\scanpage{169}
3) Le groupe $G$ des éléments inversibles de $U$
tels que \uuline{$\Delta s=s\otimes s$,} \inserted{$\varepsilon(s)=1$. Ces $s$} \deleted{\unsure{Ils}} correspondent aux
\deleted{\unsure{V}}automorphismes du foncteur $I^{\underline U}$ tels que,
si $A$ et $B$ \deleted{sont} correspondent à $V,W\in K^{\underline U}$,
alors $A\otimes B$ correspond à $V\otimes W$. [Ils
\margin{La représentation de $G$ dans $V\otimes W$ est le produit tensoriel de ses représentations dans $V$ et $W$}
% slanted
correspondent aussi aux homomorphismes
\deleted{dans $\underline U$} continus, compatibles avec $\Delta$,
dans $\underline U$, de l'algèbre $\underline U(\mathrm Z)$ correspondant
au groupe additif des entiers et les
représentations \deleted{de \unsure{groupes}}\inserted{linéaires} de dim finie
de ce groupe]. On conclut de
\margin{$G$ opère \unsure{trivialemt} dans le module trivial. La représentation de \deleted{$G$ opérant sur $V$} $G$ dans $V'$ est contragréd. de la représentation de $G$ dans $V$}
% slanted
$\Delta s=s\otimes s$ que $\varepsilon(s)=\varepsilon(s)^2$, donc $\varepsilon(s)$ est égal
à $0$ ou $1$, mais pas $0$, car $s$ est supposé
inversible, d'où \uline{$s\in G\Longrightarrow\varepsilon(s)=1$}.

De plus, la \deleted{\unsure{de}} condition $P(\unsure{w|\vee|v}\otimes1)\Delta=1$
prouve $\breve s s=1$, d'où \uline{$\breve s=s^{-1}$}.

Conclusion\,: \uline{Si $G$ est total dans $U$}, il
ne peut exister qu'une seule application
\deleted{\unsure{continue} $\unsure{\Rightarrow}$} linéaire $\vee:U\longrightarrow U$, continue,
telle que $P(\vee\otimes1)\Delta=1$. Cette $\vee$ est
automatiquement un anti-automorphisme
d'algèbre associative, satisfait à $\varepsilon(\breve x)=\varepsilon(x)$
et est \deleted{compatible avec la structure}\inserted{un automorphisme de Co-algèbre}.
Elle est aussi caractérisée par $s^\vee=s^{-1}$
si $s\in G$. \deleted{\unsure{D'où}}

$\Delta s=s\otimes s$ signifie que $f\longrightarrow\langle s,f\rangle$
\deleted{\unsure{continue}} est un hom. de l'algèbre
commutative $A=\underline U'$ (dual

% archive sheet number: 168
